%%% BEGIN fig:power-borel-plane
\begin{figure}[htbp]
\centering
\begin{tikzpicture}[fd,x=1.3cm,y=1cm,
  fd2 pole/.pic={\draw[line width=1.05pt,line cap=round] (-0.09,-0.09)--(0.09,0.09) (-0.09,0.09)--(0.09,-0.09);},
  fd2 ax/.style={lcaxis,gray!60!black},
  fd2 contour/.style={draw=#1,line width=1.0pt,postaction={decorate,decoration={markings,
      mark=at position 0.5 with {\arrow{Stealth[length=5.5pt,width=4.6pt]}}}}},
  fd2 pointer/.style={draw=#1,line width=0.45pt,-{Stealth[length=3.5pt,width=3pt]}},
  fd2 note/.style={font=\footnotesize},
  fd2 lab/.style={font=\small}]
  % axes
  \draw[fd2 ax] (-3.75,0) -- (6.05,0) node[right] {$\Real u$};
  \draw[fd2 ax] (0,-0.75) -- (0,1.75) node[left] {$\Imag u$};
  \foreach \u in {-3,-2,-1,1,2,3,4,5} \draw[gray!60!black,line width=0.4pt] (\u,-0.05) -- (\u,0.05);
  \foreach \u in {-3,-2,-1,1,2,3,4,5} \node[fd2 note] at (\u,-0.5) {$\u$};
  \node[fd2 note,anchor=north east] at (-0.03,-0.04) {$0$};
  % Borel integration paths above and below the positive axis
  \draw[fd2 contour=blue!60!black] (0,0) .. controls (0.3,0.22) and (0.55,0.24) .. (0.9,0.24) -- (5.5,0.24);
  \draw[fd2 contour=blue!60!black] (0,0) .. controls (0.3,-0.22) and (0.55,-0.24) .. (0.9,-0.24) -- (5.5,-0.24);
  \node[fd2 lab,text=blue!50!black,anchor=west] at (5.5,0.28) {$C_+$};
  \node[fd2 lab,text=blue!50!black,anchor=west] at (5.5,-0.28) {$C_-$};
  % IR renormalons (u = 2,3,...), UV renormalons (u = -1,-2,...), no pole at u = 1
  \foreach \u in {2,3,4,5} \pic[red!70!black] at (\u,0) {fd2 pole};
  \foreach \u in {-1,-2,-3} \pic[gray!55!black] at (\u,0) {fd2 pole};
  \draw[line width=0.9pt,fill=white] (1,0) circle[radius=0.085];
  % annotations
  \node[fd2 note,text=red!60!black] at (3.5,0.72) {红外重整化子：$u=2,3,\dots$};
  \node[fd2 note,align=center,text=gray!40!black] at (-1.95,0.82) {紫外重整化子：$u=-1,-2,\dots$\\（符号交错，不在积分路径上）};
  \draw[fd2 pointer=gray!60!black] (0.78,-1.02) -- (0.97,-0.12);
  \node[fd2 note,anchor=north east] at (0.95,-0.98) {无 $u=1$ 极点：不存在量纲为 2 的规范不变算符};
  \draw[fd2 pointer=red!60!black] (2.12,-1.0) -- (2.02,-0.12);
  \node[fd2 note,align=left,anchor=north west,text=red!50!black] at (1.75,-0.96) {$u=2$：沿 $C_+$ 与 $C_-$ 积分之差 $\propto(\Lambda^2/Q^2)^2$，\\对应 OPE 中的凝聚项 $\langle\frac{\alpha_s}{\pi}G^2\rangle/Q^4$};
  \node[fd2 lab] at (1.15,-2.7) {$D(Q^2)=\displaystyle\int_0^\infty\!\dd t\,\ee^{-t/\alpha_s(Q)}B[D](t)$，\quad $u=\dfrac{\beta_0t}{4\pi}$，\quad $\ee^{-t/\alpha_s(Q)}=\Big(\dfrac{\Lambda^2}{Q^2}\Big)^{u}$};
\end{tikzpicture}
\caption{大 $\beta_0$（泡链）近似下 Adler 函数 $D(Q^2)$ 的 Borel 平面（\cref{prop:power-adler-large-beta0}，部分 [P]；Borel 变换按\cref{def:power-borel}：$R=\sum_nr_n\alpha_s^{n+1}$，$B[R](t)=\sum_nr_nt^n/n!$，$u=\beta_0t/(4\pi)$，单圈跑动下 $\ee^{-t/\alpha_s(Q)}=(\Lambda^2/Q^2)^u$）。
红外重整化子极点位于正实轴 $u=2,3,\dots$（红叉），使沿正实轴的 Borel 积分不确定：绕到极点上方（$C_+$）或下方（$C_-$）的两种积分之差在 $u=u_0$ 处正比于 $(\Lambda^2/Q^2)^{u_0}$。
$u=1$ 处没有极点（空心圆），这与不存在量纲为 2 的规范不变算符相一致；$u=2$ 的极点与 OPE 中的胶子凝聚项 $\langle\frac{\alpha_s}{\pi}G^2\rangle/Q^4$ 相对应。紫外重整化子位于 $u=-1,-2,\dots$（灰叉）。
只有大 $\beta_0$ 极限中的这一结构是确立的；红外重整化子与高扭度之间的一般对应是猜想（\cref{conj:power-renormalon-ht}，[C]）。}
\label{fig:power-borel-plane}
\end{figure}
%%% END fig:power-borel-plane

%%% BEGIN fig:renormalon-chain
\begin{figure}[htbp]
\centering
\begin{tikzpicture}[fd,
  t3 bubble/.pic={\draw[quark=0.55] (-0.25,0) arc[start angle=180,end angle=0,radius=0.25];
                  \draw[quark=0.55] (0.25,0) arc[start angle=0,end angle=-180,radius=0.25];
                  \node[t3 sv] at (-0.25,0) {}; \node[t3 sv] at (0.25,0) {};},
  t3 sv/.style={circle,fill,draw=none,inner sep=0pt,minimum size=2.8pt},
  t3 bracket/.style={draw=gray!55!black,line width=0.5pt},
  t3 lab/.style={font=\small},
  t3 note/.style={font=\footnotesize}]
  % ============ (a) one-loop vertex with a dressed gluon ============
  \begin{scope}
    \coordinate (V) at (0,1.7);  \coordinate (Pin) at (-3.4,-1.0);  \coordinate (Pout) at (3.4,-1.0);
    \coordinate (A) at (-2.38,-0.19);  \coordinate (B) at (2.38,-0.19);
    \draw[photon] (0,2.95) -- (V);
    \draw[quark] (Pin) -- (A);  \draw[quark] (A) -- (V);  \draw[quark] (V) -- (B);  \draw[quark] (B) -- (Pout);
    \draw[gluon] (A) -- (-1.83,-0.19);  \pic at (-1.58,-0.19) {t3 bubble};
    \draw[gluon] (-1.33,-0.19) -- (-0.78,-0.19);  \pic at (-0.53,-0.19) {t3 bubble};
    \node[font=\small] at (0,-0.19) {$\cdots$};
    \pic at (0.53,-0.19) {t3 bubble};  \draw[gluon] (0.78,-0.19) -- (1.33,-0.19);
    \pic at (1.58,-0.19) {t3 bubble};  \draw[gluon] (1.83,-0.19) -- (B);
    \foreach \p in {V,A,B} \node[vertex] at (\p) {};
    \draw[t3 bracket] (-1.83,-0.6) -- (-1.83,-0.7) -- (1.83,-0.7) -- (1.83,-0.6);
    \node[t3 note,below=1pt] at (0,-0.7) {$n$ 个费米子泡};
    \node[t3 lab,right=1pt] at (0,2.75) {$\gamma^*$};
    \path[mom={$k$},mom sep=0.3cm,mom pos=0.5] (A) -- (-1.83,-0.19);
    \node[t3 lab] at (0,-1.65) {(a)};
  \end{scope}
  % ============ (b) Adler function: vacuum polarization with a dressed gluon ============
  \begin{scope}[xshift=7.6cm]
    \coordinate (L) at (-1.7,0); \coordinate (R) at (1.7,0);
    \coordinate (T) at (0,1.55); \coordinate (Bm) at (0,-1.55);
    \draw[photon] (-3.0,0) -- (L);  \draw[photon] (R) -- (3.0,0);
    \draw[quark] (L) arc[start angle=180,end angle=90,x radius=1.7,y radius=1.55];
    \draw[quark] (T) arc[start angle=90,end angle=0,x radius=1.7,y radius=1.55];
    \draw[quark] (R) arc[start angle=0,end angle=-90,x radius=1.7,y radius=1.55];
    \draw[quark] (Bm) arc[start angle=-90,end angle=-180,x radius=1.7,y radius=1.55];
    \draw[gluon] (T) -- (0,1.1);  \pic[rotate=90] at (0,0.85) {t3 bubble};
    \draw[gluon] (0,0.6) -- (0,0.22);
    \node[font=\small] at (0,0) {$\vdots$};
    \draw[gluon] (0,-0.22) -- (0,-0.6);  \pic[rotate=90] at (0,-0.85) {t3 bubble};
    \draw[gluon] (0,-1.1) -- (Bm);
    \foreach \p in {L,R,T,Bm} \node[vertex] at (\p) {};
    \node[t3 lab,above=1pt] at (-2.45,0.08) {$j^\mu$};
    \node[t3 lab,above=1pt] at (2.45,0.08) {$j^\nu$};
    \path[mom'={$q$},mom sep=0.3cm,mom length=0.45cm] (-3.0,0) -- (L);
    \node[t3 note,anchor=west] at (0.25,0) {$n$ 个泡};
    \node[t3 lab] at (0,-2.1) {(b)};
  \end{scope}
\end{tikzpicture}
\caption{重整化子链图（大 $\beta_0$ 近似）。胶子传播子插入 $n$ 个费米子泡：(a)~夸克--光子顶点的单圈修正；(b)~Adler 函数（流--流关联函数）中的相应图（另有两个自能型图，未画出）。
在 $\MSbar$ 方案中每个泡给出因子 $\frac{\alpha_sT_Fn_f}{3\pi}\big[\ln\frac{-k^2}{\mu^2}-\frac53\big]$（相对于裸胶子传播子）；
保留 $n_f$ 的最高次幂后作替换 $\frac43T_Fn_f\to-\beta_0$（“朴素非阿贝尔化”）并对 $n$ 求和，相当于把胶子两端的耦合换成以 $-k^2\ee^{-5/3}$ 为标度的单圈跑动耦合
$\alpha_s(\mu)\big/\big[1+\frac{\beta_0\alpha_s(\mu)}{4\pi}\ln\frac{-k^2\ee^{-5/3}}{\mu^2}\big]$。
固定 $n$ 时，圈积分中 $|k^2|\ll Q^2$ 的区域给出随 $n$ 阶乘增长的系数 $\sim n!\,\big(\frac{\beta_0}{4\pi p}\big)^n\alpha_s^{n+1}$（红外重整化子，$p=1,2,\dots$），相应 Borel 积分的不确定性 $\propto(\Lambda^2/Q^2)^p$，与高扭度贡献的量级相同；
$|k^2|\gg Q^2$ 的区域给出符号交错的紫外重整化子（见正文）。}
\label{fig:renormalon-chain}
\end{figure}
%%% END fig:renormalon-chain

%%% BEGIN fig:twist3-ssa
\begin{figure}[htbp]
\centering
\begin{tikzpicture}[fd,
  t3 hadron/.style={draw=black!70,line width=1.5pt,line cap=butt},
  t3 pole/.style={draw=red!70!black,line width=1.6pt,fd arrow=#1},
  t3 pole/.default=0.5,
  t3 lab/.style={font=\small},
  t3 note/.style={font=\footnotesize}]
  \node[jet blob,ellipse,minimum width=5.6cm,minimum height=1.1cm,inner sep=0pt] (A) at (0,-2.2) {};
  \node[jetB blob,ellipse,minimum width=5.6cm,minimum height=1.1cm,inner sep=0pt] (B) at (0,2.3) {};
  \node[hard blob,circle,minimum size=0.84cm,inner sep=0pt] (HL) at (-2.55,0) {};
  \node[hard blob,circle,minimum size=0.84cm,inner sep=0pt] (HR) at (2.55,0) {};
  \coordinate (G) at (-0.7,0);
  % hadrons
  \draw[t3 hadron] (-4.75,-2.2) -- (A.west);  \draw[t3 hadron] (A.east) -- (4.75,-2.2);
  \draw[t3 hadron] (-4.75, 2.3) -- (B.west);  \draw[t3 hadron] (B.east) -- (4.75, 2.3);
  % partons from B (gluons) and from A (quarks x1, x2)
  \draw[gluon] (-1.75,1.87) -- (HL.105);   \draw[gluon] (HR.75) -- (1.75,1.87);
  \draw[quark] (-1.75,-1.77) -- (HL.285);   \draw[quark] (HR.255) -- (1.75,-1.77);
  % soft gluon of the twist-3 correlator, attached to the final-state quark c
  \draw[gluon] (-0.7,-1.67) -- (G);
  % final-state quark c (pole propagator in red) and the observed photon
  \draw[t3 pole=0.55] (HL.east) -- (G);
  \draw[quark=0.5] (G) -- (HR.west);
  \draw[photon] (HL.55) .. controls (-1.3,1.2) and (1.3,1.2) .. (HR.125);
  \node[vertex] at (G) {};
  \draw[cut] (0,-2.95) -- (0,3.05);
  % labels
  \node[t3 lab] at (HL) {$H$};  \node[t3 lab] at (HR) {$H^\dagger$};
  \node[t3 lab] at (-1.3,-2.25) {$A^{\uparrow}$};  \node[t3 lab] at (-1.3,2.3) {$B$};
  \node[t3 lab,below=1pt] at (-4.3,-2.25) {$P_A,S_{\mathrm T}$};
  \node[t3 lab,above=1pt] at (-4.3,2.35) {$P_B$};
  \path[mom={$x_1P_A$},mom sep=0.27cm,mom length=0.45cm] (-1.75,-1.77) -- (HL.285);
  \path[mom={$x_2P_A$},mom sep=0.27cm,mom length=0.45cm] (HR.255) -- (1.75,-1.77);
  \path[mom={$x_gP_A$},mom sep=0.27cm,mom length=0.45cm,mom pos=0.45] (-0.7,-1.67) -- (G);
  \path[mom={$p_c$},mom sep=0.25cm,mom length=0.4cm] (G) -- (0,0);
  \node[t3 lab] at (0.62,1.22) {$\gamma$};
  \node[t3 note,text=red!65!black] at (-1.42,0.25) {极点};
  \node[t3 note,anchor=north] at (0,-2.95) {$T_F(x_1,x_2)$：夸克--胶子--夸克关联（扭度 3）};
\end{tikzpicture}
\caption{Qiu--Sterman 软胶子极点机制（以 $A^{\uparrow}+B\to\gamma+X$ 中的 Compton 子过程 $qg\to\gamma q$ 为例；竖直虚线为末态割线）。
横向极化强子 $A^{\uparrow}$ 的扭度 3 关联函数 $T_F(x_1,x_2)$ 提供割线两侧的夸克线（$x_1P_A$、$x_2P_A$）与一条胶子线（$x_gP_A$，$x_g\equiv x_2-x_1$），后者与末态夸克 $c$ 相连。
夹在硬部分与该胶子之间的传播子（红色）携带 $p_c-(x_2-x_1)P_A$，其分母 $\big(p_c-(x_2-x_1)P_A\big)^2+\ii0=-2p_c\cdot P_A\,(x_2-x_1-\ii0)$ 在 $x_2\to x_1$ 时为零：
$\frac{1}{x_2-x_1-\ii0}$ 的 $\ii\pi\delta(x_2-x_1)$ 部分提供单自旋不对称所需的相位，使结果正比于 $T_F(x,x)$（“软胶子极点”）。
若胶子与初态部分子相连，$\ii0$ 的符号相反，色因子也不同——这与 Sivers 函数的过程依赖性相对应（参见\cref{fig:sidis-dy-links}）。}
\label{fig:twist3-ssa}
\end{figure}
%%% END fig:twist3-ssa
